\documentclass[a4paper,12pt]{article}
\usepackage[utf8]{inputenc}
\usepackage[ngerman]{babel}
\usepackage{amsmath,amssymb}
\DeclareMathOperator{\cof}{Cof} % Cofaktormatrix
\DeclareMathOperator{\adj}{Adj} % Adjunkte Matrix
\usepackage{graphicx}
\usepackage{enumitem}
\usepackage[most]{tcolorbox} % lädt tcolorbox mit allen Libraries
\definecolor{violet}{RGB}{100, 0, 180}
\definecolor{yellow}{RGB}{5, 95, 150}
\usepackage{float}
\usepackage{bm}
\usepackage{siunitx}
\usepackage{booktabs}
\sisetup{per-mode = symbol} % statt s^{-1} bekommst du "/s"

\allowdisplaybreaks
\setlength{\parindent}{0pt}
\setlist[itemize]{noitemsep}

% Ich will, dass nach einem paragraph ein Zeilenumbruch stattfindet. Das macht dieses makro. (GPT)
\makeatletter
\renewcommand\paragraph{\@startsection{paragraph}{4}{\z@}%
	{1.5ex \@plus .5ex \@minus .2ex}% Abstand davor
	{1.0ex \@plus .2ex}%  << positiv => neue Zeile/Abstand danach
	{\normalfont\normalsize\bfseries}}% Stil
\makeatother

%% Eigene t-colorboxen
% Eigene Umgebung für Angabe-Box
\newtcolorbox{antwortbox}{
	colback=white,
	colframe=green!60!black,
	title=Antwort
}

% Eigene Umgebung für Angabe-Box
\newtcolorbox{angabebox}{
	colback=white,
	colframe=black,
	title=Angabe
}

% Eigene Umgebung für Formelsammlung-Box
\newtcolorbox{formelbox}{
	colback=white,
	colframe=orange,
	title=Formelsammlung
}

% Eigene Umgebung für Merken-Box
\newtcolorbox{merkebox}{
	colback=white,
	colframe=violet,
	title=Merken
}

% Eigene Umgebung für TODO-Box
\newtcolorbox{todobox}{
	colback=white,
	colframe=blue,
	title=TODO
}

% kompakter „File-Tag“ wie \fcolorbox{gray}{lightgray}{\texttt{…}}
\newcommand{\filebox}[1]{%
	\begingroup
	\setlength{\fboxrule}{0.6pt}% Rahmenstärke
	\setlength{\fboxsep}{1.5pt}% Innenabstand
	\fcolorbox{gray}{lightgray}{\texttt{#1}}%
	\endgroup
}
% Nutzung: \filebox{Inverse\_Matrix.md}

% Hantelzeichen für Transformationen
\newcommand{\hantel}{\mathrel{\circ\!\!-\!\!\bullet}}     % links leer, rechts gefüllt
\newcommand{\hantelrev}{\mathrel{\bullet\!\!-\!\!\circ}}  % links gefüllt, rechts leer

\title{Ökonomie der Energienetze\\\large Homework - 7}
\author{Roland Stenzl}
% \date{\today}

\begin{document}
\maketitle

\vspace{0.5cm}

\section*{Herausforderung der Amortisierung von Investitionen:
	Wasserstoff-Hochlauf und finanziellen Risiken}


\begin{angabebox}
	In Europa wird in einigen Ländern der Betrieb von Methan- und Wasserstoffnetzen durch einen gemeinsamen Betreiber angestrebt. Für eine solche Situation werden im Folgenden die Herausforderung und das Risiko der Amortisation der Investitionen in die Wasserstoffinfrastruktur für einen anschaulichen und stark vereinfachten Fall untersucht. Gehen Sie grundsätzlich bei den nachfolgenden Berechnungen davon aus, dass die beiden Netze getrennt einer sogenannten Erlösobergrenze (\textit{Revenue Cap}) unterliegen. Dabei sind folgende Annahmen zu treffen:
	
	\begin{itemize}
		\item Der aktuelle Buchwert der Vermögenswerte (d.\,h., für das Methannetz im Jahr 2023, kein Wasserstoffnetz vorhanden) beträgt 1~Mrd.~Euro;
		\item Die Erlösobergrenze für das aufzubauende Wasserstoffnetz setzt sich vereinfacht aus 80\,\% der Kapitalkosten (die restlichen 20\,\% sind für Ersatzinvestitionen im Methannetz vorgesehen) und 100\,\% der jährlichen Abschreibungen des Methannetzes zusammen.
	\end{itemize}
	
	Annahme einer linearen Abschreibung über 20 Jahre und einer Kapitalverzinsung von 6\,\% (d.\,h.\ WACC) über den gesamten Betrachtungszeitraum (2023 bis 2032). Weitere begründete Annahmen sind gegebenenfalls zu treffen.
\end{angabebox}

\subsection*{Unterpunkt a)}

\begin{angabebox}
	Berechnen Sie den sich jährlich einstellenden Endkundentarif für das Wasserstoffnetz im Betrachtungszeitraum 2023 bis 2032 unter der Annahme, dass die gebuchte Transportkapazität bis 2030 4~TWh pro Jahr beträgt. Ab dem Jahr 2030 wird von einer Verdopplung der gebuchten Transportkapazität bis zum Ende des Betrachtungszeitraums im Jahr 2032 ausgegangen. Berechnen Sie ebenfalls den über den gesamten Betrachtungszeitraum durchschnittlichen Endkundentarif.
\end{angabebox}

\paragraph{1) Jährliche Abschreibung des Methannetzes.}
Bei linearer Abschreibung über 20 Jahre mit einem Buchwert des Methannetzes von BV$_{CH_4} =$ 1 Mrd. EUR ergibt sich:

\[
\text{Dep}_{\mathrm{CH_4}} = \frac{BV_{\mathrm{CH_4}}}{20 \, a}
= \frac{1{,}000~\text{Mio.~EUR}}{20 \, a}
= 50~\text{Mio.~EUR/a}.
\]

\paragraph{2) Jährliche Kapitalkosten}
Die jährliche Kapitalverzinsung beträgt:

\[
K_{\mathrm{Zins}} = \mathrm{WACC} \cdot BV_{\mathrm{CH_4}}
\]

\paragraph{3) Erlösobergrenze des Wasserstoffnetzes.}
Die Erlösobergrenze ergibt sich zu:

\[
RC_{\mathrm{H_2}} =
0{,}8 \cdot K_{\mathrm{Zins}}
+ \text{Dep}_{\mathrm{CH_4}}
\]

\paragraph{4) Jährlich anfallender Endkundentarif.}
Der Endkundentarif ergibt sich als Quotient aus Erlösobergrenze und gebuchter Transportmenge:

\[
\tau_t = \frac{RC_{\mathrm{H_2}}}{Q_t}
\]

\begin{table}[H]
	\centering
	\label{tab:h2_tarife}
	\begin{tabular}{c c c c c c}
		\toprule
		Jahr &
		Buchwert &
		Kapitalkosten &
		Transport- &
		Erlös-OG &
		Tarif \\
		& [Mio.\,EUR] & [Mio.\,EUR] & Kapazität [TWh] & [Mio.\,EUR] & [EUR/MWh] \\
		\midrule
		2023 & 1000 & 60 & 4 & 98{,}0 & 24{,}5 \\
		2024 &  950 & 57 & 4 & 95{,}6 & 23{,}9 \\
		2025 &  900 & 54 & 4 & 93{,}2 & 23{,}3 \\
		2026 &  850 & 51 & 4 & 90{,}8 & 22{,}7 \\
		2027 &  800 & 48 & 4 & 88{,}4 & 22{,}1 \\
		2028 &  750 & 45 & 4 & 86{,}0 & 21{,}5 \\
		2029 &  700 & 42 & 4 & 83{,}6 & 20{,}9 \\
		2030 &  650 & 39 & 8 & 81{,}2 & 10{,}15 \\
		2031 &  600 & 36 & 8 & 78{,}8 & 9{,}85 \\
		2032 &  550 & 33 & 8 & 76{,}4 & 9{,}55 \\
		\bottomrule
	\end{tabular}
\end{table}

\paragraph{5) Durchschnittlicher Endkundentarif}
Für den durchschnittlichen Endkundentarif bilden wir den Mittelwert:

\[
\overline{\tau} = \frac{\sum_t \tau_t}{10} = 18{,}845 ~\text{~EUR/MWh}
\]

\subsection*{Unterpunkt b)}
	
\begin{angabebox}
	Vergleichen Sie auf Basis Ihrer Berechnungen im vorherigen Punkt die jährlichen Endkundentarife mit dem durchschnittlichen Endkundentarif. Gehen Sie dabei insbesondere auf die Vor- bzw.\ Nachteile ein, die sich aus der unterschiedlichen Tarifgestaltung ergeben können.
\end{angabebox}

\begin{figure}[h]
	\centering
	\includegraphics[width=0.8\textwidth]{Unterpunkt_b.png}
	\caption{Beschreibung der Abbildung}
	\label{fig:fig_b}
\end{figure}


In Bezug auf die \textit{willingness-to-pay} der Endkunden kann es insbesondere in der
Anfangsphase des Netzausbaus sinnvoll sein, einen durchschnittlichen Endkundentarif
anzuwenden. Bis zum Jahr 2029 liegen die jährlichen Endkundentarife über dem
durchschnittlichen Tarif, wodurch ein geglätteter Tarif zu geringeren Einstiegskosten
führt. Ab dem Jahr 2030 sinken die jährlichen Tarife infolge der erhöhten
Transportkapazität deutlich unter den durchschnittlichen Tarif.

\paragraph{Jährliche Endkundentarife.}
\textit{Vorteile:} Sinkende Tarife nach der Kapazitätserhöhung ab 2030 sowie bessere
Ansprache preissensitiver Kunden in späteren Jahren.
\textit{Nachteile:} Hohe Tarife in der Anfangsphase können den frühzeitigen
Netzanschluss und damit den Nachfragehochlauf hemmen.

\paragraph{Durchschnittlicher Endkundentarif.}
\textit{Vorteile:} Geringere Einstiegskosten erhöhen den Anreiz für einen frühen
Netzanschluss und können die Transportmengen steigern.
\textit{Nachteile:} In späteren Jahren liegt der durchschnittliche Tarif über dem
jeweils jährlichen Tarif.
In Bezug auf die \textit{willingness-to-pay} der Endkunden kann es insbesondere in der
Anfangsphase des Netzausbaus sinnvoll sein, einen durchschnittlichen Endkundentarif
anzuwenden. Bis zum Jahr 2029 liegen die jährlichen Endkundentarife über dem
durchschnittlichen Tarif, wodurch ein geglätteter Tarif zu geringeren Einstiegskosten
führt. Ab dem Jahr 2030 sinken die jährlichen Tarife infolge der erhöhten
Transportkapazität deutlich unter den durchschnittlichen Tarif.

\paragraph{Jährliche Endkundentarife.}
\textit{Vorteile:} Sinkende Tarife nach der Kapazitätserhöhung ab 2030 sowie bessere
Ansprache preissensitiver Kunden in späteren Jahren.  
\textit{Nachteile:} Hohe Tarife in der Anfangsphase können den frühzeitigen
Netzanschluss und damit den Nachfragehochlauf hemmen.

\paragraph{Durchschnittlicher Endkundentarif.}
\textit{Vorteile:} Geringere Einstiegskosten erhöhen den Anreiz für einen frühen
Netzanschluss und können die Transportmengen steigern.  
\textit{Nachteile:} In späteren Jahren liegt der durchschnittliche Tarif über dem
jeweils jährlichen Tarif.

	
\subsection*{Unterpunkt c)}
	
\begin{angabebox}
	Beschreiben Sie qualitativ den Zusammenhang zwischen dem sich einstellenden Endkundentarif und der Zahlungsbereitschaft für die Netznutzung (\textit{willingness-to-pay}). Machen Sie konkrete Vorschläge, wie Netzbetreiber die Zahlungsbereitschaft von Endkunden abschätzen können. Unterscheiden Sie dabei insbesondere zwischen Haushalts- und Industriekunden. Welche Rolle spielt dabei das Konzept der ``Opportunitätskosten''?
\end{angabebox}

\end{document}documentclass[a4paper,12pt]{article}
\usepackage[utf8]{inputenc}
\usepackage[ngerman]{babel}
\usepackage{amsmath,amssymb}
\DeclareMathOperator{\cof}{Cof} % Cofaktormatrix
\DeclareMathOperator{\adj}{Adj} % Adjunkte Matrix
\usepackage{graphicx}
\usepackage{enumitem}
\usepackage[most]{tcolorbox} % lädt tcolorbox mit allen Libraries
\definecolor{violet}{RGB}{100, 0, 180}
\definecolor{yellow}{RGB}{5, 95, 150}
\usepackage{float}
\usepackage{bm}
\usepackage{siunitx}
\usepackage{booktabs}
\sisetup{per-mode = symbol} % statt s^{-1} bekommst du "/s"

\allowdisplaybreaks
\setlength{\parindent}{0pt}
\setlist[itemize]{noitemsep}

% Ich will, dass nach einem paragraph ein Zeilenumbruch stattfindet. Das macht dieses makro. (GPT)
\makeatletter
\renewcommand\paragraph{\@startsection{paragraph}{4}{\z@}%
	{1.5ex \@plus .5ex \@minus .2ex}% Abstand davor
	{1.0ex \@plus .2ex}%  << positiv => neue Zeile/Abstand danach
	{\normalfont\normalsize\bfseries}}% Stil
\makeatother

%% Eigene t-colorboxen
% Eigene Umgebung für Angabe-Box
\newtcolorbox{antwortbox}{
	colback=white,
	colframe=green!60!black,
	title=Antwort
}

% Eigene Umgebung für Angabe-Box
\newtcolorbox{angabebox}{
	colback=white,
	colframe=black,
	title=Angabe
}

% Eigene Umgebung für Formelsammlung-Box
\newtcolorbox{formelbox}{
	colback=white,
	colframe=orange,
	title=Formelsammlung
}

% Eigene Umgebung für Merken-Box
\newtcolorbox{merkebox}{
	colback=white,
	colframe=violet,
	title=Merken
}

% Eigene Umgebung für TODO-Box
\newtcolorbox{todobox}{
	colback=white,
	colframe=blue,
	title=TODO
}

% kompakter „File-Tag“ wie \fcolorbox{gray}{lightgray}{\texttt{…}}
\newcommand{\filebox}[1]{%
	\begingroup
	\setlength{\fboxrule}{0.6pt}% Rahmenstärke
	\setlength{\fboxsep}{1.5pt}% Innenabstand
	\fcolorbox{gray}{lightgray}{\texttt{#1}}%
	\endgroup
}
% Nutzung: \filebox{Inverse\_Matrix.md}

% Hantelzeichen für Transformationen
\newcommand{\hantel}{\mathrel{\circ\!\!-\!\!\bullet}}     % links leer, rechts gefüllt
\newcommand{\hantelrev}{\mathrel{\bullet\!\!-\!\!\circ}}  % links gefüllt, rechts leer

\title{Ökonomie der Energienetze\\\large Homework - 7}
\author{Roland Stenzl}
% \date{\today}

\begin{document}
\maketitle

\vspace{0.5cm}

\section*{Herausforderung der Amortisierung von Investitionen:
	Wasserstoff-Hochlauf und finanziellen Risiken}


\begin{angabebox}
	In Europa wird in einigen Ländern der Betrieb von Methan- und Wasserstoffnetzen durch einen gemeinsamen Betreiber angestrebt. Für eine solche Situation werden im Folgenden die Herausforderung und das Risiko der Amortisation der Investitionen in die Wasserstoffinfrastruktur für einen anschaulichen und stark vereinfachten Fall untersucht. Gehen Sie grundsätzlich bei den nachfolgenden Berechnungen davon aus, dass die beiden Netze getrennt einer sogenannten Erlösobergrenze (\textit{Revenue Cap}) unterliegen. Dabei sind folgende Annahmen zu treffen:
	
	\begin{itemize}
		\item Der aktuelle Buchwert der Vermögenswerte (d.\,h., für das Methannetz im Jahr 2023, kein Wasserstoffnetz vorhanden) beträgt 1~Mrd.~Euro;
		\item Die Erlösobergrenze für das aufzubauende Wasserstoffnetz setzt sich vereinfacht aus 80\,\% der Kapitalkosten (die restlichen 20\,\% sind für Ersatzinvestitionen im Methannetz vorgesehen) und 100\,\% der jährlichen Abschreibungen des Methannetzes zusammen.
	\end{itemize}
	
	Annahme einer linearen Abschreibung über 20 Jahre und einer Kapitalverzinsung von 6\,\% (d.\,h.\ WACC) über den gesamten Betrachtungszeitraum (2023 bis 2032). Weitere begründete Annahmen sind gegebenenfalls zu treffen.
\end{angabebox}

\subsection*{Unterpunkt a)}

\begin{angabebox}
	Berechnen Sie den sich jährlich einstellenden Endkundentarif für das Wasserstoffnetz im Betrachtungszeitraum 2023 bis 2032 unter der Annahme, dass die gebuchte Transportkapazität bis 2030 4~TWh pro Jahr beträgt. Ab dem Jahr 2030 wird von einer Verdopplung der gebuchten Transportkapazität bis zum Ende des Betrachtungszeitraums im Jahr 2032 ausgegangen. Berechnen Sie ebenfalls den über den gesamten Betrachtungszeitraum durchschnittlichen Endkundentarif.
\end{angabebox}

\paragraph{1) Jährliche Abschreibung des Methannetzes.}
Bei linearer Abschreibung über 20 Jahre mit einem Buchwert des Methannetzes von BV$_{CH_4} =$ 1 Mrd. EUR ergibt sich:

\[
\text{Dep}_{\mathrm{CH_4}} = \frac{BV_{\mathrm{CH_4}}}{20 \, a}
= \frac{1{,}000~\text{Mio.~EUR}}{20 \, a}
= 50~\text{Mio.~EUR/a}.
\]

\paragraph{2) Jährliche Kapitalkosten}
Die jährliche Kapitalverzinsung beträgt:

\[
K_{\mathrm{Zins}} = \mathrm{WACC} \cdot BV_{\mathrm{CH_4}}
\]

\paragraph{3) Erlösobergrenze des Wasserstoffnetzes.}
Die Erlösobergrenze ergibt sich zu:

\[
RC_{\mathrm{H_2}} =
0{,}8 \cdot K_{\mathrm{Zins}}
+ \text{Dep}_{\mathrm{CH_4}}
\]

\paragraph{4) Jährlich anfallender Endkundentarif.}
Der Endkundentarif ergibt sich als Quotient aus Erlösobergrenze und gebuchter Transportmenge:

\[
\tau_t = \frac{RC_{\mathrm{H_2}}}{Q_t}
\]

\begin{table}[H]
	\centering
	\label{tab:h2_tarife}
	\begin{tabular}{c c c c c c}
		\toprule
		Jahr &
		Buchwert &
		Kapitalkosten &
		Transport- &
		Erlös-OG &
		Tarif \\
		& [Mio.\,EUR] & [Mio.\,EUR] & Kapazität [TWh] & [Mio.\,EUR] & [EUR/MWh] \\
		\midrule
		2023 & 1000 & 60 & 4 & 98{,}0 & 24{,}5 \\
		2024 &  950 & 57 & 4 & 95{,}6 & 23{,}9 \\
		2025 &  900 & 54 & 4 & 93{,}2 & 23{,}3 \\
		2026 &  850 & 51 & 4 & 90{,}8 & 22{,}7 \\
		2027 &  800 & 48 & 4 & 88{,}4 & 22{,}1 \\
		2028 &  750 & 45 & 4 & 86{,}0 & 21{,}5 \\
		2029 &  700 & 42 & 4 & 83{,}6 & 20{,}9 \\
		2030 &  650 & 39 & 8 & 81{,}2 & 10{,}15 \\
		2031 &  600 & 36 & 8 & 78{,}8 & 9{,}85 \\
		2032 &  550 & 33 & 8 & 76{,}4 & 9{,}55 \\
		\bottomrule
	\end{tabular}
\end{table}

\paragraph{5) Durchschnittlicher Endkundentarif}
Für den durchschnittlichen Endkundentarif bilden wir den Mittelwert:

\[
\overline{\tau} = \frac{\sum_t \tau_t}{10} = 18{,}845 ~\text{~EUR/MWh}
\]

\subsection*{Unterpunkt b)}
	
\begin{angabebox}
	Vergleichen Sie auf Basis Ihrer Berechnungen im vorherigen Punkt die jährlichen Endkundentarife mit dem durchschnittlichen Endkundentarif. Gehen Sie dabei insbesondere auf die Vor- bzw.\ Nachteile ein, die sich aus der unterschiedlichen Tarifgestaltung ergeben können.
\end{angabebox}

\begin{figure}[h]
	\centering
	\includegraphics[width=0.8\textwidth]{Unterpunkt_b.png}
	\caption{Beschreibung der Abbildung}
	\label{fig:fig_b}
\end{figure}


In Bezug auf die \textit{willingness-to-pay} der Endkunden kann es insbesondere in der
Anfangsphase des Netzausbaus sinnvoll sein, einen durchschnittlichen Endkundentarif
anzuwenden. Bis zum Jahr 2029 liegen die jährlichen Endkundentarife über dem
durchschnittlichen Tarif, wodurch ein geglätteter Tarif zu geringeren Einstiegskosten
führt. Ab dem Jahr 2030 sinken die jährlichen Tarife infolge der erhöhten
Transportkapazität deutlich unter den durchschnittlichen Tarif.

\paragraph{Jährliche Endkundentarife.}
\textit{Vorteile:} Sinkende Tarife nach der Kapazitätserhöhung ab 2030 sowie bessere
Ansprache preissensitiver Kunden in späteren Jahren.
\textit{Nachteile:} Hohe Tarife in der Anfangsphase können den frühzeitigen
Netzanschluss und damit den Nachfragehochlauf hemmen.

\paragraph{Durchschnittlicher Endkundentarif.}
\textit{Vorteile:} Geringere Einstiegskosten erhöhen den Anreiz für einen frühen
Netzanschluss und können die Transportmengen steigern.
\textit{Nachteile:} In späteren Jahren liegt der durchschnittliche Tarif über dem
jeweils jährlichen Tarif.
In Bezug auf die \textit{willingness-to-pay} der Endkunden kann es insbesondere in der
Anfangsphase des Netzausbaus sinnvoll sein, einen durchschnittlichen Endkundentarif
anzuwenden. Bis zum Jahr 2029 liegen die jährlichen Endkundentarife über dem
durchschnittlichen Tarif, wodurch ein geglätteter Tarif zu geringeren Einstiegskosten
führt. Ab dem Jahr 2030 sinken die jährlichen Tarife infolge der erhöhten
Transportkapazität deutlich unter den durchschnittlichen Tarif.

\paragraph{Jährliche Endkundentarife.}
\textit{Vorteile:} Sinkende Tarife nach der Kapazitätserhöhung ab 2030 sowie bessere
Ansprache preissensitiver Kunden in späteren Jahren.  
\textit{Nachteile:} Hohe Tarife in der Anfangsphase können den frühzeitigen
Netzanschluss und damit den Nachfragehochlauf hemmen.

\paragraph{Durchschnittlicher Endkundentarif.}
\textit{Vorteile:} Geringere Einstiegskosten erhöhen den Anreiz für einen frühen
Netzanschluss und können die Transportmengen steigern.  
\textit{Nachteile:} In späteren Jahren liegt der durchschnittliche Tarif über dem
jeweils jährlichen Tarif.

	
\subsection*{Unterpunkt c)}
	
\begin{angabebox}
	Beschreiben Sie qualitativ den Zusammenhang zwischen dem sich einstellenden Endkundentarif und der Zahlungsbereitschaft für die Netznutzung (\textit{willingness-to-pay}). Machen Sie konkrete Vorschläge, wie Netzbetreiber die Zahlungsbereitschaft von Endkunden abschätzen können. Unterscheiden Sie dabei insbesondere zwischen Haushalts- und Industriekunden. Welche Rolle spielt dabei das Konzept der ``Opportunitätskosten''?
\end{angabebox}

\end{document}
